\section{The model}
\label{sec:model}

\paragraph{Effort profiles.}
A team of $n$ agents produces an \emph{effort profile} $S \in \cube$, where
$S_i = 1$ means agent~$i$ exerted effort. We take $S$ uniform on $\cube$; the
consequences of that choice are discussed in Section~\ref{sec:limitations}.

\paragraph{The record.}
A \emph{record} of size $K$ is a stochastic encoder
$\enc : \cube \to \simplex(\labels)$, written $\enc[s,y] = \Prob(Y = y \mid S = s)$.
The realised label $Y \in \labels$ is everything that persists. Both uses below
read $Y$ and nothing else.

\paragraph{Use one: reconstruction.}
A reconstruction rule $\rec : \labels \to \cube$ recovers an estimate of the
profile. Its cost is expected Hamming error,
\[
  D(\enc,\rec) \;=\; \E\bigl[\dH\bigl(S, \rec(Y)\bigr)\bigr],
\]
and $\Dnull(K) = \min_{\enc,\rec} D(\enc,\rec)$ is the unconstrained optimum.
$\Dnull(K)$ is the Hamming-cube $K$-median value and is classical; nothing in
this paper claims it is new.

\paragraph{Use two: deferred payment.}
Each agent~$i$ is paid by a bounded decoder $\pay_i \in [0,1]^K$ that also sees
only $Y$. Write the \emph{unilateral deviation vector} of agent~$i$ at teammate
profile $\sminus \in \{0,1\}^{n-1}$ as
\[
  d_i(\sminus)[y] \;=\; \enc\bigl[(1,\sminus),\,y\bigr] - \enc\bigl[(0,\sminus),\,y\bigr],
\]
the change in the label distribution caused by agent~$i$ alone. The record
\emph{implements margin $\gamma$} if there are bounded decoders with
\[
  \bigl\langle \pay_i,\, d_i(\sminus) \bigr\rangle \;\ge\; \gamma
  \qquad \text{for every } i \text{ and every } \sminus .
\tag{M}\label{eq:margin}
\]

\paragraph{What \eqref{eq:margin} is, and what it is not.}
We call $\gamma$ the \emph{accountability margin}, or equivalently the payment
margin, and we deliberately avoid calling \eqref{eq:margin} incentive
compatibility. It says the record lets a bounded payment reward unilateral effort
by at least $\gamma$; it does not say effort is a best response, because effort
\emph{costs} are not modelled. A complete strategic statement would fix a cost
$c_i$ and ask for $\gamma \ge c_i$, at which point \eqref{eq:margin} becomes the
information-side half of an incentive constraint. That is the sense in which the
results below are a prerequisite for mechanism design rather than an instance of
it: they bound what \emph{any} bounded payment can extract from the record, so a
record failing \eqref{eq:margin} at margin $\gamma$ cannot support effort of cost
$\gamma$ under any payment rule, whatever the equilibrium concept. This is not a
full strategic equilibrium model, and no claim below should be read as one.

Condition~\eqref{eq:margin} is also deliberately strong in a second way: it is an
\emph{all-teammate-profile} requirement, not an equilibrium condition. Exerting
effort must pay, whatever the rest of the team did. This makes our impossibility
results stronger and our constructions harder, which is the conservative
direction for both.

\paragraph{Labels, not bits.}
$K$ counts \emph{labels}---how many episodes the surviving record can tell
apart. A fixed-length encoding of that record costs $\log_2 K$ bits. The two are
kept separate throughout: thresholds and frontiers are stated in $K$, because
that is what the combinatorics is about, and the engineering readings in
Section~\ref{sec:agents} convert to bits explicitly. A statement like ``$K$ must
roughly double'' is one retained bit, not a doubling of the memory bill.

\paragraph{What may be random, and what may persist.}
$\enc$ is an arbitrary stochastic map, so the encoder may flip coins freely; that
generality is what makes the impossibility results strong, since they hold
against randomised records too. Two things are excluded and both matter.
The coins are \emph{not} retained: only the realised label $Y$ persists, so the
decoders cannot condition on the randomness that produced it. And no auxiliary
state escapes the count---Remark~\ref{rem:counting} is the formal statement.
Without those two restrictions accountability is free and vacuous: a record that
privately stores one bit per agent and reports a single label would satisfy
\eqref{eq:margin} at $K=1$. Encoder-side randomness therefore cannot buy margin
by itself; what it can do is spread a profile's mass across labels, which is
exactly what the construction of Lemma~\ref{lem:sufficiency} uses.

\paragraph{The object of study.}
\[
  \Dopt(K,\gamma) \;=\; \min\bigl\{\, D(\enc,\rec) \;:\; (\enc,\rec)
  \text{ implements margin } \gamma \,\bigr\}.
\]

\begin{remark}[All decoder-side state counts toward $K$]
\label{rem:counting}
$K$ counts everything the decoders may condition on. Without this the tradeoff
is vacuous: a seed, a receipt, an agent identifier or any auxiliary random
variable kept outside the counted record restores unlimited sensitivity at no
reconstruction cost. The theorems below are statements about what a
\emph{self-contained} record of size $K$ can do.
\end{remark}

\begin{table}[t]
\centering
\caption{The model, read as an agent system.}
\label{tab:mapping}
\small
\begin{tabular}{@{}ll@{}}
\toprule
formal object & agent-system reading \\
\midrule
$S_i$        & whether agent $i$ did its assigned work \\
$S$          & the joint contribution profile of the episode \\
$\enc$       & the compaction step: summarisation, eviction, sampling \\
$Y$          & the surviving record: what a later reader actually has \\
$K$          & how many episodes the record can tell apart \\
$\rec$       & post-hoc reconstruction of who contributed \\
$\pay_i$     & deferred credit, payment or blame assigned to agent $i$ \\
$\gamma$     & how strongly the record must reward unilateral effort \\
$\dH$        & error in the reconstructed contribution profile \\
\bottomrule
\end{tabular}
\end{table}

\begin{definition}[Flatness threshold]
\label{def:kflat}
$\Kflat(n)$ is the least $K$ at which accountability is free, i.e.\ at which
$\Dopt(K,\gamma) = \Dnull(K)$ for some $\gamma > 0$.
\end{definition}
